\section{四交比 corner の完全判定と六点配置による十分性}
\label{sec:six-point-classification}
本節では，前節までに証明した十一の必要条件が十分でもあることを示す．
その鍵は，四交比トーラスを\emph{標識付き可換三次曲面の六点配置空間}と同定することである．
以下，基礎体は標数0の代数閉体とし，十八の係数$w_{1bc},w_{2bc}$はすべて非零とする．

\begin{theorem}[corner有限性の必要十分条件]
\label{thm:general-position-iff}
$B(w)=\operatorname{Jac}(\widetilde Q,W(w))$，$e=e_{A_0}$，
$C(w)=B(w)/B(w)eB(w)$と置く．
$Q=(q_{bc})$，$q_{bc}=w_{2bc}/w_{1bc}$とすると，次は同値である．
\begin{enumerate}
\item $\dim_kC(w)<\infty$．
\item $Q$の九つの$2\times2$小行列式がすべて非零であり，
$\det Q\ne0$かつ$\det Q^{\langle-1\rangle}\ne0$である．
\item $P=(Q^{\langle-1\rangle})^{\mathsf t}$の三つの列$p_1,p_2,p_3$と，
座標点$p_4=(1:0:0),p_5=(0:1:0),p_6=(0:0:1)$が，
$\mathbb P^2$の一般位置の六点をなす．
\end{enumerate}
従って，四交比トーラスにおける有限corner locusは，
\emph{正確に}前節の十一のdivisorの補集合である．
この開集合は，六点の順序を保つ標識付き可換三次曲面のモジュライ
$M_{\mathrm{cubic}}$と自然に同型になる．
\end{theorem}
\begin{proof}
(i)$\Rightarrow$(ii)は定理\ref{thm:torus}と
命題\ref{prop:two-determinants}による．
(ii)$\Leftrightarrow$(iii)を示す．$P$の各成分は$q_{bc}^{-1}\ne0$であり，
その$2\times2$小行列式は，添字と非零係数を除けば
$Q$の対応する小行列式と等しい．さらに
$\det P=\det Q^{\langle-1\rangle}$である．従って(ii)のうち
九小行列式と後者の行列式は，順に「可動点二つと座標点一つが非共線」
「可動点三つが非共線」を表す．非零成分は「可動点一つと座標点二つが非共線」を表す．
座標点三つはもとより非共線である．

残る条件は六点が一つの二次曲線上にないことである．
座標三点を通るconicの方程式は
$\alpha X_1X_2+\beta X_2X_3+\gamma X_3X_1=0$
であり，可動点$p_i=(P_{0i}:P_{1i}:P_{2i})$への評価行列は
\[
 M_{i\bullet}=(P_{0i}P_{1i},P_{1i}P_{2i},P_{2i}P_{0i})
 =(P_{0i}P_{1i}P_{2i})
 (P_{2i}^{-1},P_{0i}^{-1},P_{1i}^{-1}).
\]
すべての因子が非零なので$\det M\ne0$は
$\det\bigl((P^{\langle-1\rangle})^{\mathsf t}\bigr)=\det Q\ne0$
と同値である．これで(ii)$\Leftrightarrow$(iii)を得る．

最後に(iii)$\Rightarrow$(i)を証明する．
$S=\operatorname{Bl}_{p_1,\ldots,p_6}\mathbb P^2$は滑らかな可換三次曲面で，
$H=-K_S$はampleである．AOU\cite[\S\S4,8，特に(4.1)，(4.7)，(4.9)，(8.4)--(8.7)]{AOU}
の九線束$E_v$とそのrolled-up helix代数を取る．
その幾何学的potentialの二つの関係系列は，基底を適切に規格化すると
\begin{equation}\label{eq:geometric-relations}
 w^{\mathrm{geo}}_{1bc}=P_{cb},\qquad
 w^{\mathrm{geo}}_{2bc}=1\qquad(0\le b,c\le2)
\end{equation}
となる．これは次の初等的な式からも直接確認できる．
$u_b=(P_{0b}P_{1b},P_{1b}P_{2b},P_{2b}P_{0b})$，
$L_{b,c}(X)=\det[p_{b+1},p_{c+4},X]$とし，
$D_b(X)=(u_{b+1}\times u_{b+2})\cdot
(X_1X_2,X_2X_3,X_3X_1)$（添字を$3$で読む）とする．
$D_b$は$p_{b+1}$以外の五点を通るconicである．
また$M_b(X)=\det[p_{b+2},p_{b+3},X]$と書くと，
行列式の展開により全ての$c$について
\begin{equation}\label{eq:six-point-identities}
 \sum_{b=0}^2P_{cb}D_b(X)L_{b,c}(X)=0,\qquad
 \sum_{b=0}^2M_b(X)L_{b,c}(X)=0
\end{equation}
が成立する．各関係空間はAOU(4.2)により一次元であるため，
\eqref{eq:six-point-identities}が\eqref{eq:geometric-relations}を与える．
これらの多項式恒等式は独立の整数係数による厳密検算コードにも収めた．

\eqref{eq:geometric-relations}より，その二系列の比は
$w^{\mathrm{geo}}_{2bc}/w^{\mathrm{geo}}_{1bc}
=P_{cb}^{-1}=q_{bc}$である．
元の二系列と幾何学的二系列は各矢$y_{bc}$の非零rescalingで一致する．
残りの$w_{0bc}$は頂点$A_0$を除いた商$C$に影響しないから，
\[
 C(w)\simeq C\bigl(w^{\mathrm{geo}}\bigr)
\]
が道長次数付き代数として成り立つ．

幾何学的代数については，AOU(4.9)により
\begin{equation}\label{eq:geometric-rolled-orbit}
 B\bigl(w^{\mathrm{geo}}\bigr)
 \simeq\bigoplus_{v,v'}\bigoplus_{n\ge0}
 \operatorname{Hom}_S(E_v,E_{v'}(nH))
\end{equation}
がcut次数付き代数として成立し，$e=e_{0,0}$に対応する線束は
$E_{0,0}=\mathcal O_S$である．
任意の$v$について$m\gg0$を固定すれば
$E_v^\vee(mH)$は大域生成される．評価写像の核を$K_v$として
\[
 0\longrightarrow K_v\longrightarrow
 H^0(S,E_v^\vee(mH))\otimes\mathcal O_S
 \longrightarrow E_v^\vee(mH)\longrightarrow0
\]
を取る．これを$E_{v'}((n-m)H)$でtensorし，
Serre消滅$H^1(S,K_v\otimes E_{v'}((n-m)H))=0$（$n\gg0$）を使うと，
乗法写像
\begin{equation}\label{eq:serre-factorization}
 H^0(S,E_v^\vee(mH))\otimes
 H^0(S,E_{v'}((n-m)H))
 \twoheadrightarrow
 \operatorname{Hom}_S(E_v,E_{v'}(nH))
\end{equation}
が全射となる．左辺の積は必ず
$E_v\to\mathcal O_S(mH)\to E_{v'}(nH)$という，
構造頂点$e$を通る経路に対応する．頂点は有限個なので全ての$v,v'$に共通の
十分大きい$N$を取れば，
$B_n\subset BeB$（$n\ge N$）が従う．
cut次数零は有限次元で，各斉次成分も有限次元であるから，
$\dim_kC\bigl(w^{\mathrm{geo}}\bigr)<\infty$となる．
これで(i)が証明された．

最後に，標識付き六点のmoduliで座標三点$p_4,p_5,p_6$を固定すると，
残る自由度は$P$の各列のスカラーと座標三点を保つ対角$\mathrm{PGL}_3$で割ることである．
この商はちょうど$Q$の行・列rescaling商であり，
$P=(Q^{\langle-1\rangle})^{\mathsf t}$およびその逆は
全成分非零のトーラス上で正則だから，主張した開集合の同型が得られる．
\end{proof}

\begin{corollary}[AS heart と anticanonical heart の一致領域]
\label{cor:complete-as-ac}
十八係数非零の係数チャートで，標準AOU型のfolding，
反Serre作用のcut shiftとの同定，必要なnoetherian性を満たすAS対象に限る．
このとき自然なcorner restrictionが
\[
 \mathcal A_{\mathrm{AS}}\xrightarrow{\sim}
 \mathcal A_{\mathrm{ac}}=\qgr R(\mathcal O,\tau)
\]
という構造対象・反Serre作用と両立するabel同値であることと，
定理\ref{thm:general-position-iff}の十一条件は同値である．
\end{corollary}
\begin{proof}
定理\ref{thm:general-position-iff}と命題\ref{prop:corner}，
式\eqref{eq:orbit}を適用する．
\end{proof}

\begin{remark}
定理\ref{thm:general-position-iff}は，\emph{元の27係数potential全体が}
可換三次曲面から生じると主張するものではない．
$w_{0bc}$は任意でよく，幾何学的潜在係数と一致させる必要はない．
同定しているのは$A_0$除去後のcorner商である．
また，十一条件に違反した場合に任意の別のabel同値・導来同値が存在しないという主張でもない．
Rains heart，VdB heartとの比較は依然として独立の課題である．
\end{remark}
